% !TeX root = ../main.tex
\section{Определение позиции на открытой местности}
Для начала выберем место около которого не будет зданий в окрестности 50 метров. Для этого было выбрано место обзорной площадки в ландшафтом парке Митино. 

\begin{figure}[H]
    \centering
    \includegraphics[width=0.5\textwidth]{5274195380109254549.jpg}
    \caption{Местоположение обзорной площадки}
\end{figure}

Прождав 5 минут, находясь на площадке с включенным GPS, получили результат, представленный ниже.

\begin{figure}[H]
    \centering
    \includegraphics[width=0.5\textwidth]{5274195380109254551.jpg}
    \caption{Оценка местоположения}
\end{figure}

Как можно заметить, положение точки истинного положения и определённого с помощью спутника отличается примерно на 14-15 метров. Ожидается что данный результат погрешности будет только расти в последующих опытах из-за перекрытия обзоров спутноков.
Также пронаблюдаем колличесвто спутников, которые смог поймать приемник на момент фиксации результата.

\begin{figure}[H]
    \centering
    \includegraphics[width=0.5\textwidth]{спутники 1.jpg}
    \caption{Колличество спутников пойманных в момент фиксации}
\end{figure}

В данном опыте было обнаружено 11 спутников потенциально способствующих определению местоположения.

\begin{figure}[H]
    \centering
    \includegraphics[width=0.5\textwidth]{5274195380109254553.jpg}
    \caption{Фотоотчет о выполнении}
\end{figure}